\documentclass[11pt]{article}

% ----------------------------------------------------------------------------
% E-HANK: baseline model note (self-contained).
% Documents ONLY the configuration actually run: core (domestic-good) numeraire,
% import booking of the green margin, with the current calibration. Switched-off
% ARS features (energy in production, subsistence, GHH, DCP export pricing,
% importer weights in the fund return, monetary inertia / wage-inflation
% response, pass-through to the green price) and the alternative numeraire /
% booking are omitted. epsT (untargeted lump sum) and ishock (monetary shock)
% are retained as hooks even though zero in the baseline experiments.
% ----------------------------------------------------------------------------

\usepackage[margin=1in]{geometry}
\usepackage{amsmath,amssymb,mathtools}
\usepackage{booktabs}
\usepackage{enumitem}
\usepackage[hidelinks]{hyperref}

\newcommand{\I}{\mathbb{I}}
\newcommand{\E}{\mathbb{E}}
\newcommand{\num}{\mathrm{num}}
\newcommand{\stst}{\mathrm{ss}}
\newcommand{\dsw}{D^{\mathrm{sw}}}
\newcommand{\DG}{D^{G}}
\newcommand{\PEB}{P^{E}_{B}}
\newcommand{\PEG}{P^{E}_{G}}
\newcommand{\pE}{p^{E}}
\newcommand{\coh}{\mathrm{coh}}

\title{\textbf{E-HANK: baseline model and equations}\\[2pt]
\large A small open economy HANK with a brown/green durable adoption margin}
\author{Guillaume Schwegler, Boris Chafwehe, Francesca Diluiso}
\date{Working note --- \today}

\begin{document}
\maketitle

\begin{abstract}
\noindent This note states the baseline model implemented in the code, with all
equilibrium equations. The aggregate structure is \cite{ARS2024} (henceforth
ARS): a small open economy with sticky wages, imported energy and foreign goods
with sticky retail prices, and a mutual fund holding claims on domestic firms,
importers and a share of the energy endowment. The only structural addition is
a \emph{discrete brown/green durable adoption choice} in the household block;
the aggregate blocks follow ARS's quantitative model (including foreign-owned
importers). The note documents the baseline
configuration only: domestic-good numeraire and import booking of the green
margin. Sections~\ref{sec:hh}--\ref{sec:agg} give the household problem and its
aggregation; Sections~\ref{sec:firms}--\ref{sec:clearing} give firms, prices,
the open economy, policy and market clearing; Section~\ref{sec:ss} gives the
steady state, shocks and calibration.
\end{abstract}

% ============================================================================
\section{Conventions and notation}
% ============================================================================

Time is quarterly. A subscript $t$ is suppressed where unambiguous; $x(-1)$ and
$x(+1)$ denote lag and (expected) lead. A bar denotes a steady-state value,
e.g.\ $\bar\pE$. All \emph{relative prices} carry the subscript ``\,$/P$\,''
where $P$ is the CPI, e.g.\ $\PEB\equiv \PEB{}_{,t}/P_t$; these are the objects
the CES demand system uses.

\paragraph{Unit of account.} The unit of account is the domestic good. The
household reads, solves in units of account, and returns assets converted to CPI
units by $A^{CPI}=A\,p^{\num}$, with
\begin{equation}
p^{\num}=p_{H/P},\qquad
1+r^{\num}=(1+r)\,\frac{p_{H/P}(-1)}{p_{H/P}},\qquad
atw^{\num}=\frac{atw}{p_{H/P}} .
\end{equation}
Objects contractually denominated in the CPI---fiscal transfers and the
switching cost $\psi_g$---are divided by $p^{\num}$ on the way into the budget.
Taking the domestic good rather than the CPI as unit of account makes the
resource constraint independent of how the CPI weights brown and green energy.

% ============================================================================
\section{Households}\label{sec:hh}
% ============================================================================

There are three permanent discount-factor types $i\in\{0,1,2\}$ of equal mass,
with $\beta_i$ (Section~\ref{sec:cal}); within a type, households face
idiosyncratic labour productivity and hold a discrete durable technology.

\subsection{Idiosyncratic productivity and the durable state}

Labour productivity $e$ follows a Rouwenhorst discretisation of an AR(1) in
logs, $\log e' = \rho_e \log e + \varepsilon$, with $n_e$ points, persistence
$\rho_e$ and innovation s.d.\ $\sigma_e$; transition matrix $\Pi$, mean-one grid.

Each household carries two binary durable indicators: the technology it
\emph{operates this period} $g\in\{B,G\}$ and the technology it \emph{held last
period} $g_-\in\{B,G\}$. We collapse the pair into a four-point index
\begin{equation}
d\in\{\,BB,\ BG,\ GB,\ GG\,\},\qquad
\text{(now, before)}=\{(B,B),(B,G),(G,B),(G,G)\}.
\end{equation}
Define the indicators
\begin{equation}
\mathrm{IS\_GREEN}(d)=\I[g=G]=(0,0,1,1),\qquad
\mathrm{PAYS\_SWITCH}(d)=\I[g=G,\,g_-=B]=(0,0,1,0).
\end{equation}
The switching cost $\psi_g$ is incurred exactly in the transition $g_-=B\to g=G$
(state $GB$). Brown is \emph{absorbing}; a green durable \emph{breaks down} to
brown at rate $\delta_g$. The breakdown Markov (applied in the \texttt{dep}
stage, rows = state this period, columns = state next period) is
\begin{equation}
M_{\mathrm{dep}}=
\begin{pmatrix}
1 & 0 & 0 & 0\\
1 & 0 & 0 & 0\\
\delta_g & 0 & 0 & 1-\delta_g\\
\delta_g & 0 & 0 & 1-\delta_g
\end{pmatrix}.
\end{equation}
On the equilibrium path $BG$ carries zero mass (green is lost only through
breakdown, recorded as $g_-=B$), so the effective state is three-point. In
steady state a permanent green share is sustained by a flow of switchers,
\begin{equation}
\dsw=\delta_g\,\DG. \label{eq:stockflow}
\end{equation}

Forward stage order within a period is
$\texttt{[dep, prod, durables, consav]}$: breakdown, productivity, the discrete
adoption choice, and the EGM consumption--savings step. The adoption stage uses
the discrete--continuous structure of \cite{DCEGM2017}.

\subsection{Preferences and the type-specific bundle price}

Flow felicity is CRRA over a consumption bundle $c$ with disutility of labour
handled at the union level (Section~\ref{sec:union}):
\begin{equation}
u(c)=\begin{cases}\log c & \text{eis}=1\\[2pt]
\dfrac{c^{1-1/\text{eis}}}{1-1/\text{eis}} & \text{eis}\neq 1.\end{cases}
\end{equation}
The bundle $c$ is a CES nest of energy against a home/foreign composite, with
ARS parameters $(\alpha_E,\eta_E)$. Because a green adopter buys energy at $\PEG$
and a brown household at $\PEB$, each durable type faces its own energy price and
hence its own \emph{exact} cost-of-living index:
\begin{align}
\pE_d &= \PEB + \mathrm{IS\_GREEN}(d)\,\bigl(\PEG-\PEB\bigr),\\
p_d &= \Bigl[\,1+\alpha_E\bigl((\pE_d)^{1-\eta_E}-(\PEB)^{1-\eta_E}\bigr)\Bigr]^{\frac{1}{1-\eta_E}}
\quad(\eta_E\neq 1),\qquad
p_d=\bigl(\pE_d/\PEB\bigr)^{\alpha_E}\ (\eta_E=1).
\label{eq:pd}
\end{align}
Equation~\eqref{eq:pd} uses the equilibrium CES identity
$\alpha_E(\PEB)^{1-\eta_E}+(1-\alpha_E)p_{HF/P}^{1-\eta_E}=1$
(the block \texttt{CESprices}, eq.~\eqref{eq:innernest}), so \emph{brown}
households have $p_d\equiv 1$ at every date (exact nesting). $p_d$ is
CPI-relative; the budget uses $p_d^{\num}\equiv p_d/p^{\num}$.

\subsection{Cash on hand}\label{sec:coh}

Let $a$ be beginning-of-period assets (in units of account, on a grid with lower
bound $a\ge \underline a$, $\underline a=0$). Cash on hand of a household in
durable state $d$ with productivity $e$ is
\begin{equation}
\coh(d,e,a)=(1+r^{\num})\,a+atw^{\num}\,e
+\frac{T^{\text{fisc}}}{p^{\num}}
-\frac{\psi_g}{p^{\num}}\,\mathrm{PAYS\_SWITCH}(d),
\label{eq:coh}
\end{equation}
with the (CPI-denominated) transfer
\begin{equation}
T^{\text{fisc}} = \epsilon^{T}+\iota\,(\pE-\bar\pE)\,\bar c^{B}_{E,i}(e,a)
\qquad\text{(untargeted lump sum $\epsilon^{T}$ + Slutsky transfer).}
\end{equation}
Here $\iota=\texttt{insE}$ scales the Slutsky transfer, $\epsilon^{T}=\texttt{epsT}$
is an untargeted lump sum, and $\bar c^{B}_{E,i}(e,a)$ is household $i$'s
steady-state \emph{brown} energy demand, collapsed over the durable axis and
broadcast identically across durable states (so a switcher keeps the same
transfer; the switching margin is undistorted). Both instruments are zero in the
baseline. The switching cost enters cash on hand only in state $GB$.

\subsection{Discrete adoption choice}\label{sec:logit}

In the \texttt{durables} stage the household chooses its current technology $g$
given the technology it holds (post-breakdown), by a logit over continuation
values with taste scale $\sigma_\varepsilon=\texttt{taste\_shock}$. Feasible
transitions carry the flow payoff
\begin{equation}
f(d\mid d_-)=
\begin{cases}
0 & \text{stay/renew brown: } (BB\mid BB),(BB\mid BG),\\
-\,\texttt{green\_block}\ \ (\text{masked to }-\infty\text{ if }\coh(GB)\le 0)
& \text{switch to green: } (GB\mid BB),(GB\mid BG),\\
0 & \text{stay/renew green: } (GG\mid GB),(GG\mid GG),\\
-\infty & \text{otherwise (infeasible).}
\end{cases}
\end{equation}
The switching cost itself is already in $\coh(GB)$ through
eq.~\eqref{eq:coh}. \texttt{green\_block} is a counterfactual wedge, $0$ in the
baseline (margin open) and $+\infty$ (large) in the ``no-adoption'' variant
(margin shut, used for the adoption-channel decomposition). Given post-choice
values $V(d)$, the choice probabilities and the expected (pre-shock) value are
the standard logit / log-sum-exp:
\begin{equation}
\Pr(d\mid d_-)=\frac{\exp\!\bigl([f(d\mid d_-)+V(d)]/\sigma_\varepsilon\bigr)}
{\sum_{d'}\exp\!\bigl([f(d'\mid d_-)+V(d')]/\sigma_\varepsilon\bigr)},
\qquad
\tilde V(d_-)=\sigma_\varepsilon\log\!\sum_{d'}\exp\!\Bigl(\tfrac{f(d'\mid d_-)+V(d')}{\sigma_\varepsilon}\Bigr).
\end{equation}

\subsection{Consumption--savings (EGM)}

Conditional on the durable state, the household solves a standard
consumption--savings problem with the type-specific bundle price $p_d^{\num}$.
The budget, Euler and envelope are
\begin{align}
p_d^{\num}\,c + a' &= \coh(d), & a'&\ge \underline a, \label{eq:budget}\\
u'(c) &= \beta_i\,p_d^{\num}\,\E\,V_a(a',e',d') & &\text{(Euler)},\\
V_a &= (1+r^{\num})\,u'(c)/p_d^{\num} & &\text{(envelope)},\\
V &= u(c)+\beta_i\,\E\,V(a',e',d') & &\text{(value)},
\end{align}
solved by the endogenous grid method, with $u'(c)=c^{-1/\text{eis}}$.

\subsection{Household outputs (hetoutputs)}\label{sec:hetout}

Given the policy $c(d,e,a)$, the household's \emph{exact} CES demand split by
durable type is
\begin{align}
c_E(d) &= \alpha_E\bigl(\pE_d/p_d\bigr)^{-\eta_E} c, &
c_{HF}(d) &= (1-\alpha_E)\bigl(p_{HF/P}/p_d\bigr)^{-\eta_E} c, \label{eq:demand}\\
c_E^{B}(d) &= c_E(d)\,\bigl(1-\mathrm{IS\_GREEN}(d)\bigr), &
c_E^{G}(d) &= c_E(d)\,\mathrm{IS\_GREEN}(d).
\end{align}
(Prices in \eqref{eq:demand} are CPI-relative; the denominator is $p_d$, not
$p_d^{\num}$.) The block also returns $d_G(d)=\mathrm{IS\_GREEN}(d)$,
$d_{sw}(d)=\mathrm{PAYS\_SWITCH}(d)$, the MPC, and flow utility $u(c)$.

% ============================================================================
\section{Aggregation}\label{sec:agg}
% ============================================================================

For each type $i$, SSJ aggregates the hetoutputs over the stationary
distribution $\Lambda_i$ into $C_i,A_i,\mathrm{MPC}_i,\DG_i,\dsw_i,
C^{B}_{E,i},C^{G}_{E,i},C_{HF,i},U_i$, and the model averages equally across the
three types:
\begin{equation}
X=\tfrac13\textstyle\sum_i X_i,\qquad
X\in\{C,A,\mathrm{MPC},\DG,\dsw,C^{B}_{E},C^{G}_{E},C_{HF},\ldots\},
\qquad \beta=\tfrac13\textstyle\sum_i\beta_i .
\end{equation}
Assets in units of account are converted to CPI units, $A^{CPI}=A\,p^{\num}$.
Aggregate demands for the home good, the foreign good and energy are then built
from the household aggregates (block \texttt{hh\_outputs\_dur}), applying on top
only the inner home/foreign nest, common to all households:
\begin{equation}
c_E = C^{B}_{E}+C^{G}_{E},\qquad
c_H = (1-\alpha_F)\,p_{H/HF}^{-\eta}\,C_{HF},\qquad
c_F = \alpha_F\,p_{F/HF}^{-\eta}\,C_{HF}. \label{eq:cHcF}
\end{equation}
Using $\sum_d c_E(d)$ rather than the CES demand at the aggregate index avoids a
Jensen bias of order $0.5\%$.

% ============================================================================
\section{Firms and domestic production}\label{sec:firms}
% ============================================================================

A competitive final domestic good is produced from labour. With productivity $Z$
and a constant gross markup $\mu_{\stst}$,
\begin{equation}
y = Z\,n,\qquad
\mathrm{gdp}=Z\,n,\qquad
\text{btw}_n = w\,n,\qquad
atw_n = (1-\tau^{Y})\,w\,n,
\end{equation}
where $y$ is home-good output, $n$ hours, $w$ the real wage. Monopolistic
profits are capitalised into the mutual fund:
\begin{equation}
\mathrm{div}=(\mu_{\stst}-1)\,w\,n,\qquad
J=\mathrm{div}+\frac{J(+1)}{1+r^{\text{ante}}},\qquad
j=\frac{J(+1)}{1+r^{\text{ante}}}.
\end{equation}

\subsection{Wage setting (sticky wages)}\label{sec:union}

A union sets the nominal wage subject to Calvo frictions $\theta_w$, giving a
wage Phillips curve. With $\kappa_w=(1-\theta_w)(1-\beta\theta_w)/\theta_w$,
labour-disutility shifter $\varphi$ (\texttt{vphi}), Frisch elasticity
$\texttt{frisch}$, and a real-wage-rigidity term controlled by $w_{BG}$,
\begin{align}
\pi^{w}(1+\pi^{w})
&=\beta\,\pi^{w}(+1)\bigl(1+\pi^{w}(+1)\bigr)
+\frac{\kappa_w}{\varphi\,n^{1/\text{frisch}}}\;\Xi,\\
\Xi &= \varphi\,n^{1/\text{frisch}}
-\frac{atw\,C^{-1/\text{eis}}}{\mu_{\stst}}
- w_{BG}\,\underbrace{\frac{(w-\bar w)\,w\,\bar C^{-1/\text{eis}}\,\bar n}{\mu_{\stst}\,n\,\bar w}}_{\text{real-wage rigidity}},
\end{align}
and the nominal and real wage evolve as $W=(1+\pi^{w})W(-1)$, $w=W/P$. Here
$atw=atw_n/n$ is the after-tax wage per hour.

% ============================================================================
\section{Prices}\label{sec:prices}
% ============================================================================

All retail prices are relative to the CPI $P$; $Q$ is the real exchange rate
(price of the foreign basket in home units), $P^{*}_{F},P^{E*}$ are exogenous
world prices.

\paragraph{Domestic good.}
\begin{equation}
P_H=\frac{\mu_{\stst}}{Z}\,W,\qquad
p_{H/P}=P_H/P,\qquad
\pi^{H}=P_H/P_H(-1)-1,\qquad w=W/P.
\end{equation}

\paragraph{Imported foreign good (Calvo $\theta_F$).} With
$\kappa_F=(1-\theta_F)(1-\beta_F\theta_F)/\theta_F$, $\beta_F=1/(1+\bar r^{\text{ante}})$,
\begin{equation}
\kappa_F\Bigl(\frac{Q P^{*}_{F}}{p_{F/P}}-1\Bigr)
+\beta_F\,\pi^{F}(+1)\bigl(1+\pi^{F}(+1)\bigr)
-\pi^{F}(1+\pi^{F})=0,\qquad p_{F/P}=P_F/P .
\end{equation}

\paragraph{Energy retail price and the price cap (Calvo $\theta_E$).} The
wholesale/retail brown energy price $\pE=P_E/P$ satisfies a NKPC, and the
\emph{household} brown price applies the cap:
\begin{align}
\kappa_E\Bigl(\frac{Q P^{E*}}{\pE}-1\Bigr)
+\beta_E\,\pi^{E}(+1)\bigl(1+\pi^{E}(+1)\bigr)
-\pi^{E}(1+\pi^{E})&=0,
& \kappa_E&=(1-\theta_E)(1-\beta_E\theta_E)/\theta_E,\\
\PEB &= (1-\tau^{E})\,\pE+\tau^{E}\,\bar\pE,
&& \tau^{E}=1\Rightarrow \PEB\ \text{pinned at }\bar\pE .
\label{eq:cap}
\end{align}

\paragraph{Green energy price.} With green/brown ratio $\varrho=\texttt{pE\_g\_ratio}$,
the green price is fixed relative to the pre-crisis brown price:
\begin{equation}
\PEG=\varrho\,\bar\pE . \label{eq:pEG}
\end{equation}

\paragraph{CPI (brown-anchored; ``Option~C'').} The energy leg of the CPI is
anchored on the brown price alone. With $\alpha=\alpha_E+(1-\alpha_E)\alpha_F$,
the inner (energy vs home/foreign) and outer (home vs foreign) CES identities
that define $p_{HF/P}$ and $p_{H/HF}$ are
\begin{align}
(1-\alpha_E)\,p_{HF/P}^{1-\eta_E}+\alpha_E\,(\PEB)^{1-\eta_E}&=1, \label{eq:innernest}\\
(1-\alpha_F)\,p_{H/HF}^{1-\eta}+\alpha_F\,p_{F/HF}^{1-\eta}&=1, \label{eq:outernest}
\end{align}
with $p_{F/HF}=p_{F/P}/p_{HF/P}$, $p_{H/P}=p_{H/HF}\,p_{HF/P}$, $p_H^{*}=p_{H/P}/Q$.
The gap between this measured CPI and the share-weighted (``true'') deflator is
computed \emph{ex post}, outside the model:
\begin{equation}
\phi^{1-\eta_E}=1+\alpha_E\,\DG\bigl[(\PEG)^{1-\eta_E}-(\PEB)^{1-\eta_E}\bigr].
\end{equation}
Under the domestic-good numeraire this choice affects only the measured deflator
and the monetary rule, not the resource constraint.

\paragraph{Aggregate price level and inflation.}
\begin{equation}
P=(1+\pi)P(-1),\qquad
\pi^{\text{ann}}=(1+\pi)^4-1,\ \text{etc.}
\end{equation}

% ============================================================================
\section{Open economy}\label{sec:open}
% ============================================================================

\paragraph{Foreign demand for the home good.} With foreign demand shifter
$\alpha^{*}$, trade elasticity $\gamma$, and world demand $C^{*}$,
\begin{equation}
c_H^{*}=\alpha^{*}\,\bigl(p_H^{*}\bigr)^{-\gamma}\,C^{*}.
\end{equation}

\paragraph{Uncovered interest parity.} The real exchange rate solves
\begin{equation}
\frac{Q}{Q(+1)}\,(1+r^{\text{ante}})=1+r^{*}. \label{eq:uip}
\end{equation}

\paragraph{Importer margins (foreign-owned).} Retail importers of the foreign
good and of energy are \emph{owned by foreigners} (ARS, \S4.1), so their margins
over the world price accrue abroad and are \emph{not} capitalised into domestic
wealth. They enter the model only by valuing the import bill at the retail (not
world) price in the balance of payments below:
\begin{equation}
D_F = \bigl(p_{F/P}-Q P^{*}_{F}\bigr)c_F,\qquad
D_E = \bigl(\pE-Q P^{E*}\bigr)c_E .
\end{equation}
Both margins are zero in steady state ($p_{F/P}=Q P^{*}_{F}$, $\pE=Q P^{E*}$), so
this convention leaves the steady state unchanged.

\paragraph{World energy supply (IEA block).} World supply is a fixed quantity
$E^{\text{shock}}$ plus intertemporal stock arbitrage with cost $\Gamma_{\text{arb}}$:
\begin{equation}
E_{\text{stock}} = \frac{P^{E*}(+1)/(1+r^{*})-P^{E*}}{\Gamma_{\text{arb}}},\qquad
E^{\text{supply}} = E^{\text{shock}}+\bigl(E_{\text{stock}}(-1)-E_{\text{stock}}\bigr).
\end{equation}
In the baseline the home economy owns \emph{none} of the world energy endowment
($\zeta^{E}=0$: a pure importer), so the associated rents are not capitalised
into domestic wealth. The endowment-ownership channel ($\zeta^{E}>0$, ARS
\S3.4/\S4.1) is retained in the code and used only as a robustness/energy-independence
experiment.

\paragraph{Fund return (revaluation).} The realised (ex post) return on the
fund's claims defines $r$:
\begin{equation}
1+r=\frac{J}{j(-1)}.
\end{equation}
Importers are foreign-owned and the energy endowment is not held domestically
($\zeta^{E}=0$), so the only domestically held risky claim is domestic firms:
$A^{\text{dom}}=j$, with return $r^{\text{dom}}=r$; $A^{\text{dom}}$ enters only
the current-account revaluation term below.

\paragraph{Balance of payments.}\label{sec:bop}
The green margin enters the import booking through two terms,
\begin{equation}
\text{imports}_{\text{dur}}=\psi_g\,\dsw
\quad\text{(switching cost booked as import)},\qquad
\text{energy\_gap}=(\PEB-\PEG)\,C^{G}_{E}
\quad\text{(adopters' import saving)}.
\end{equation}
Exports, imports and net exports are
\begin{align}
\text{exports}&=Q\,p_H^{*}c_H^{*},\\
\text{imports}&=p_{F/P}\,c_F+\pE\,c_E+\text{imports}_{\text{dur}}-\text{energy\_gap},\\
nx&=Q\,p_H^{*}c_H^{*}-p_{F/P}\,c_F-\pE\,c_E-\text{imports}_{\text{dur}}+\text{energy\_gap}.
\end{align}
Net foreign assets accumulate with a revaluation term:
\begin{align}
\text{reval} &= (r-r^{\text{ante}}(-1))A^{CPI}(-1)-(r^{\text{dom}}-r^{\text{ante}}(-1))A^{\text{dom}}(-1),\\
\mathrm{nfa} &= nx+\text{reval}+(1+r^{\text{ante}}(-1))\,\mathrm{nfa}(-1),
\qquad nx/\mathrm{gdp}=nx/y. \label{eq:nfa}
\end{align}

% ============================================================================
\section{Policy}\label{sec:policy}
% ============================================================================

\subsection{Monetary policy}

The nominal rate rule and the Fisher ex-ante real rate are
\begin{equation}
1+i = (1+\bar r^{*})\,(1+\phi_\pi\pi)\,(1+\phi_{\pi e}\pi(+1))+\text{ishock},
\qquad
r^{\text{ante}} = \frac{1+i}{1+\pi(+1)}-1.
\end{equation}
The baseline is the ARS \emph{constant real rate} rule
$(\phi_\pi,\phi_{\pi e})=(0,1)\Rightarrow r^{\text{ante}}=\bar r^{*}$; a standard
Taylor rule sets $(\phi_\pi,\phi_{\pi e})=(1.5,0)$. \texttt{ishock} is a monetary
policy shock (zero unless a monetary experiment is run).

\subsection{Fiscal policy}\label{sec:fiscal}

The government runs the two energy-crisis instruments and an untargeted lump sum,
all debt financed, with debt stabilised through the distortionary labour tax:
\begin{align}
\text{Subsidy} &= \tau^{E}\,(\pE-\bar\pE)\,c_E
&& \text{price cap (cost scales with actual use),}\\
\text{Ttargeted} &= \iota\,(\pE-\bar\pE)\,\bar C^{B}_{E}
&& \text{Slutsky transfer (cost scales with pre-crisis brown use),}\\
\text{Tuntargeted} &= \epsilon^{T}, &&\\[2pt]
\text{spending} &= \text{Tuntargeted}+\text{Ttargeted}+\text{Subsidy},\qquad
\text{taxation}=\tau^{Y}\,w\,n. &&
\end{align}
The government budget constraint and the debt-stabilising tax rule are
\begin{align}
B &= (1+r^{\text{ante}}(-1))\,B(-1)+\text{spending}-\text{taxation},
\label{eq:gbc}\\
\tau^{Y} &= \psi_B\bigl(B(-1)-\bar B\bigr),\qquad \psi_B=\texttt{psiB}.
\label{eq:taurule}
\end{align}
In the baseline $\bar B=0$ (no steady-state government debt): crisis outlays are
deficit-financed and repaid by raising the labour tax in proportion to the debt's
deviation from zero. Government debt $B$ is a genuine transition unknown (it moves
in the dynamics) and enters household wealth through asset-market clearing,
eq.~\eqref{eq:assets}.

% ============================================================================
\section{Market clearing and equilibrium}\label{sec:clearing}
% ============================================================================

The equilibrium residuals are
\begin{align}
\text{goods:}&\quad c_H+c_H^{*}-y=0, \label{eq:goods}\\
\text{assets (Walras):}&\quad A^{CPI}-\mathrm{nfa}-j-B=0, \label{eq:assets}\\
\text{energy:}&\quad
\begin{cases}
P^{E*}-P^{E*,\text{shock}}=0 & \varepsilon^{E}=\infty\ \text{(price taker; price shock)},\\
c_E-E^{\text{supply}}=0 & \varepsilon^{E}<\infty\ \text{(fixed quantity; supply shock).}
\end{cases}\label{eq:energy}
\end{align}
Equation~\eqref{eq:assets} is a Walras'-law residual (not itself a solved
target); it is monitored to machine precision as a consistency check.

\paragraph{Transition unknowns/targets.} The linear (SSJ) transition solves for
$\{y,\pi,r,\tau^{Y},P^{E*},w\}$ against
$\{\text{goods},\ \pi\text{ residual},\ r\text{ residual},\ \tau^{Y}\text{ residual},
\ \text{energy},\ w\text{ residual}\}$.

% ============================================================================
\section{Steady state, shocks, calibration}\label{sec:ss}
% ============================================================================

\subsection{Openness normalisations}

With $\zeta^{E}=0$ (pure importer) the ARS home-endowment rescaling is inactive,
so the normalisations reduce to
\begin{equation}
\alpha_F=\frac{\alpha-\alpha_E}{1-\alpha_E},\qquad
\gamma=\frac{\chi_{\text{tgt}}-(1-\alpha)(1-\alpha_F)\eta_E}{(1-\alpha)\alpha_F+\alpha},\quad \eta=\gamma,\ \ \chi_{\text{tgt}}=0.3,
\end{equation}
\begin{equation}
y=1,\qquad Z=\mu_{\stst}=1.03,\qquad \alpha^{*}=\alpha,\qquad E^{\text{shock}}=\alpha_E .
\end{equation}
(For the robustness runs with $\zeta^{E}>0$ the endowment rescaling
$y=1-\zeta^{E}\alpha_E$, $\mu_{\stst}\leftarrow\mu_{\stst}(1-\zeta^{E}\alpha_E)$,
$\alpha^{*}=\alpha-\zeta^{E}\alpha_E$ applies.)

\subsection{Steady-state unknowns/targets}

The steady state solves for $\{\varphi,\ \beta_{\max},\ y,\ \psi_g\}$ against
$\{\pi^{w}\text{ res},\ \mathrm{nfa}\text{ res},\ \text{goods},\ \DG-0.05\}$: $\psi_g$
is calibrated to the level target $\DG_{\stst}=0.05$ and $\beta_{\max}$ closes net
foreign assets. In the ``no-adoption'' counterfactual ($\texttt{green\_block}$
large) $\DG\equiv 0$ for every $\psi_g$, so $\psi_g$ is not identifiable and is
carried over, fixed, from the matching adoption steady state.

\subsection{Shocks}

\begin{align}
\text{Brown-price shock (requires }\varepsilon^{E}=\infty):&\quad
P^{E*,\text{shock}}_t=\text{size}\cdot\rho^{t},\quad \rho=2^{-1/\text{half-life}},\\
\text{Brown-supply shock (requires }\varepsilon^{E}<\infty):&\quad
E^{\text{shock}}_t=\bar E^{\text{shock}}-\text{drop}\cdot\bar E^{\text{shock}}\ \text{for }t<\text{quarters},\ \text{else }\bar E^{\text{shock}} .
\end{align}
The baseline price shock is $\text{size}=1$ (world price doubles on impact),
half-life $16$; the baseline supply shock is $-10\%$ for $6$ quarters.

\subsection{Calibration}\label{sec:cal}

\begin{table}[h]\centering\small
\begin{tabular}{lll}
\toprule
Parameter & Value & Source / target\\
\midrule
$\alpha_E$ energy expenditure share & 0.04 & ARS\\
$\eta_E$ energy substitution & 0.10 & ARS\\
$\alpha$ total import share & 0.30 & ARS\\
$r$, $\mu_{\stst}$ (primitive) & 0.01, 1.03 & ARS\\
eis, Frisch & 1, 0.5 & ARS\\
$\rho_e,\sigma_e,n_e$ & 0.92, 0.57, 7 & ARS\\
$\beta$ types (3), spread & 0.06 & implied MPC $\approx 0.30$\\
$\theta_w,\theta_E,\theta_F$ (Calvo) & 0.938, 0.65, 0.90 & ARS\\
$w_{BG}$ real-wage rigidity & 5 & ARS\\
$\zeta^{E}$ home energy-rent share & 0 & pure importer (baseline); $>0$ for robustness\\
\midrule
$\delta_g$ green breakdown rate & 0.05 & assumption ($\approx 5$-yr life)\\
$\varrho=\PEG/\PEB$ steady-state gap & 0.80 & heat-pump / EV delivered-cost ratio\\
$\psi_g$ switching cost & 0.538 & \emph{solved} to $\DG_{\stst}=0.05$\\
$\sigma_\varepsilon$ logit taste scale & 0.05 & calibrated jointly with $\psi_g$\\
\midrule
$\tau^{E},\iota,\epsilon^{T}$ policy & 0 & laissez-faire baseline\\
$B$ steady-state debt & 0 & no SS government debt\\
$\psi_B$ debt-feedback in $\tau^{Y}$ & 0.04 & \\
\bottomrule
\end{tabular}
\caption{Baseline calibration. $\psi_g$ and $\sigma_\varepsilon$ are jointly
calibrated to the \emph{level} target $\DG_{\stst}=0.05$ and are not separately
identified by an adoption-\emph{elasticity} moment (an open item).}
\end{table}

% ============================================================================
\begin{thebibliography}{9}
\bibitem{ARS2024} Auclert, A., Monnery, H., Rognlie, M., Straub, L. (2024).
\emph{Managing an Energy Shock}. Working paper.
\bibitem{DCEGM2017} Iskhakov, F., J{\o}rgensen, T., Rust, J., Schjerning, B.
(2017). \emph{The endogenous grid method for discrete-continuous dynamic choice
models with (or without) taste shocks}. Quantitative Economics 8(2), 317--365.
\end{thebibliography}

\end{document}
